\section{Taller resuelto: Interpretar un refinamiento realmente ejecutado}

\noindent\textbf{Nivel:} Intermedio.

\apartado{1. Enunciado}{En el ensayo libre de 60 s del T-6 se usaron pasos 0.10, 0.05 y 0.025 s. Compare las TAS finales y calcule diferencias sucesivas. ¿Permiten asignar un orden universal a JSBSim?}

\apartado{2. Datos}{Los tres CSV y registros de versión están incluidos en el proyecto; las condiciones iniciales y duración son comunes.}

\apartado{3. Incógnitas}{Las magnitudes y comprobaciones solicitadas en el caso.}

\apartado{4. Sistema de referencia}{Misma propiedad TAS SI al final de cada trayectoria.}

\apartado{5. Hipótesis físicas}{Comparación de un experimento libre concreto, sin atribuir todo el error a una sola potencia de paso.}

\apartado{6. Modelo matemático}{Interpretar un refinamiento realmente ejecutado}

\apartado{7. Ecuaciones simbólicas}{Se parte de los balances y definiciones siguientes.}

\begin{equation}\Delta_1=|V_{0.10}-V_{0.05}|\end{equation}

\begin{equation}\Delta_2=|V_{0.05}-V_{0.025}|\end{equation}

\begin{equation}p_{observado}=\ln(\Delta_1/\Delta_2)/\ln2\end{equation}

\apartado{8. Desarrollo algebraico}{Las operaciones se realizan en este orden, conservando las condiciones del modelo.}

\begin{enumerate}

\item Se extraen los valores finales sin redondear antes de restarlos.

\item Se calculan diferencias y un orden observado sólo como indicador condicionado.

\item La mezcla de actualización de masa, tolerancia de trim e integradores impide generalizarlo a todos los estados y casos.

\end{enumerate}

\apartado{9. Sustitución numérica}{Se usan las magnitudes del apartado de datos en SI; los ángulos se convierten a radianes antes de evaluar funciones trigonométricas. Las densidades indicadas como ISA se obtienen de la capa correspondiente.}

\noindent Diferencia primera:\quad $5.151754832866118 \cdot 10^{-5}\simeq 5.151755\times10^{-5}\,\mathrm{m\,s^{-1}}$.\par

\noindent Diferencia segunda:\quad $2.087958689855895 \cdot 10^{-5}\simeq 2.087959\times10^{-5}\,\mathrm{m\,s^{-1}}$.\par

\noindent Orden observado condicionado:\quad $1.302970770044629\simeq 1.302971\,$.\par

\apartado{10. Resultados con unidades}{Los resultados pertenecen al modelo y al intervalo de datos declarados.}

\begin{center}\begin{tabular}{lr}\toprule Magnitud&Resultado\\\midrule

Diferencia primera & $5.151755\times10^{-5}\,\mathrm{m\,s^{-1}}$\\

Diferencia segunda & $2.087959\times10^{-5}\,\mathrm{m\,s^{-1}}$\\

Orden observado condicionado & $1.302971\,$\\

\bottomrule\end{tabular}\end{center}

\apartado{11. Interpretación}{Las diferencias pequeñas y decrecientes aportan sensibilidad temporal para esta TAS final; no prueban un orden formal global.}

\apartado{12. Comprobación independiente y alcance}{Comprobar también altura, actitud y masa final en los registros y repetir con otro paso si la aplicación exige tolerancias más estrictas.}

\apartado{13. Limitaciones}{Las autonomías completas del 747 no se declaran convergentes por este ensayo corto de otra aeronave.}

\par\smallskip\noindent{\footnotesize\textit{Datos y solución originales. Las sustituciones registradas se reproducen con SymPy y Studium; eso no establece validación experimental ni revisión experta.}}
